\chapter{Uvod}
\label{ch:uvod}

Poglavje predstavi motivacijo za lokalno zaznavanje nevarnih
zvočnih dogodkov. Nato opredeli cilje in raziskovalna vprašanja, glavne
prispevke ter zgradbo preostalega dela.

\section{Motivacija in opredelitev problema}

Zvočni dogodki, kot so sirena, razbitje stekla, strel ali pasji
lajež, lahko človeka opozorijo na nevarnost, še preden jo opazi. Gluhe in
naglušne osebe zvočnega opozorila ne morejo vedno zaznati. Namen rešitve je zato
zvok iz okolice sproti pretvoriti v vidno informacijo na zaslonu. Cilj dela ni
razviti medicinskega ali varnostno certificiranega pripomočka, temveč preveriti,
ali je mogoče celotno pot od mikrofona do uporabniškega opozorila izvesti na
računsko in pomnilniško omejeni napravi.

Običajen pristop bi lahko zvočne posnetke pošiljal v oddaljeno računalniško
okolje, kjer bi jih obdelal večji model. Takšna rešitev je odvisna od omrežne
povezave, hkrati pa neprekinjen prenos zvoka iz domačega okolja odpira vprašanja
zasebnosti. V tej nalogi se vsi koraki izvajajo lokalno na razvojni plošči, brez
shranjevanja zvočnih podatkov. Takšna zasnova zmanjšuje izpostavljenost zasebnih
podatkov in omogoča delovanje brez dostopa do interneta.

Lokalna izvedba prinese strožje tehnične omejitve. Mikrokrmilnik mora
neprekinjeno zajemati zvok, izračunati strnjeni spektrogram, izvesti
nevronsko mrežo, filtrirati nestabilne napovedi in osveževati zaslon. Vse to
mora opraviti z omejenim pomnilnikom in dovolj hitro, da med računanjem ne
izgubi novih vzorcev. Poleg hitrosti je pomembno tudi zavračanje običajnih
zvokov. Naloga zato obravnava tako izvajanje modela z nevronskim
pospeševalnikom kot tudi merjenje pravilnih zaznav in lažnih alarmov na napravi.

\section{Cilji in raziskovalna vprašanja}

Glavni cilj je razviti delujoč vgrajeni sistem na razvojni plošči
STM32N6570-DK, ki z vgrajenim mikrofonom zaznava izbrane nevarne zvočne
dogodke in rezultat brez omrežne povezave prikaže na zaslonu. Sistem razlikuje
med naslednjimi zvočnimi razredi:

\begin{itemize}
\item pasji lajež,
\item razbitje stekla,
\item strel,
\item sirena,
\item govor in
\item ostalo.
\end{itemize}

Poleg zvočnih razredov lahko sistem prikaže še dve stanji:

\begin{itemize}
\item negotovo napoved in
\item čakanje na dogodek.
\end{itemize}

Dodatni cilj je pripraviti ponovljiv postopek vrednotenja, s katerim je mogoče
točno različico modela in programa povezati z izmerjenimi rezultati.

V nalogi odgovarjamo na naslednja raziskovalna vprašanja:

\begin{enumerate}
  \item Ali je mogoče kvantizirani model za zaznavanje zvočnih dogodkov
        izvajati v realnem času na mikrokrmilniku STM32N6 z uporabo
        nevronskega pospeševalnika Neural-ART?
  \item Kakšen delež potrjenih pravilnih zaznav in lažnih nevarnostnih alarmov
        doseže končni sistem pri nadzorovanem predvajanju zvokov na fizični
        napravi?
  \item Kako se posamezni zvočni razredi razlikujejo po uspešnosti klasifikacije in kateri
        dejavniki omejujejo delovanje sistema?
\end{enumerate}

Prvo vprašanje obravnava čas izvajanja in porabo pomnilnika, drugo uspešnost
zaznavanja in lažne alarme pri preizkusu na fizični napravi, tretje pa razlike
v uspešnosti klasifikacije med razredi in omejitve sistema.

Model in odločitvena pravila so bili izbrani z učnimi in razvojnimi podatki.
Končni preizkus je uporabil 60 novih zvočnih virov, ki pri učenju ali izbiri
pragov niso bili uporabljeni.

\section{Prispevki diplomskega dela}

Glavni prispevki dela so:

\begin{itemize}
  \item povezava neprekinjenega zajema z vgrajenega mikrofona, digitalne
        predobdelave, modela YAMNet-1024 in nevronskega pospeševalnika v
        samostojno delujočo aplikacijo;
  \item priprava ciljnega učnega postopka z izrezovanjem kratkotrajnih dogodkov,
        povečanjem podatkov ter dodatnima razredoma govor in ostalo;
  \item izvedba časovnega filtra s pragovi po razredih ter uporabniškega
        vmesnika, ki prikazuje trenutno in nedavne napovedi;
  \item lokalna zasnova brez pošiljanja zvoka v računalniški oblak in brez
        trajnega shranjevanja surovih zvočnih vzorcev v običajnem delovanju;
  \item sledljiv realni preizkus s 60 novimi, slušno pregledanimi zvočnimi
        dražljaji, shranjenimi dnevniki, kontrolnimi vsotami in ponovljivo
        analizo rezultatov.
\end{itemize}

Prednaučeni model YAMNet-1024 ni bil razvit v okviru te naloge. V nalogi smo
naučili novo izhodno plast, ki izhod modela razvrsti v izbrane zvočne razrede.
Pripravili smo tudi podatke, model prilagodili za izvajanje na mikrokrmilniku,
razvili pravila za prikaz rezultatov in ovrednotili delovanje celotnega sistema.

\section{Opis vsebine dela}

V poglavju \ref{ch:teorija} so predstavljene osnove digitalnega
zvoka, strnjenega spektrograma, nevronskih mrež in kvantizacije. Poglavje
\ref{ch:zasnova} opredeli zahteve, strojno platformo ter tok podatkov, poglavje
\ref{ch:implementacija} pa podrobno opiše zajem zvoka, predobdelavo, pripravo
modela, izvajanje na pospeševalniku in uporabniški vmesnik. Metodologija
razvojnih in končnih preskusov je podana v poglavju \ref{ch:metodologija},
izmerjeni rezultati v poglavju \ref{ch:rezultati}, odgovori na raziskovalna
vprašanja, omejitve projekta in sklepne ugotovitve pa v poglavju
\ref{ch:razprava-sklep}.
